% Artifact appendix in the CGO artifact evaluation layout
% (https://ctuning.org/ae/submission_extra.html). \input this file after
% \appendix in the paper, or build main.tex for a standalone copy.
%
% Fill in before submission: \artifactdoi (the Zenodo DOI of the archived
% branch), and the table numbers in \papertable and \paperablation.
\providecommand{\artifactdoi}{\textit{(Zenodo DOI)}}
\providecommand{\artifactrepo}{\url{https://github.com/wraith1995/Halide/tree/inductive_artifact}}
\providecommand{\papertable}{the benchmark table}
\providecommand{\paperablation}{the ablation table}

\section{Artifact Appendix}

\subsection{Abstract}

The artifact is a fork of the Halide compiler with inductive Funcs, and the
benchmark suite behind the paper's evaluation. Each benchmark builds the same
recurrence in two ways: as an inductive Func, scheduled with fusion, sliding
window and storage folding, and as an update definition over an RDom, which
materializes the whole trajectory. Both are compared against the fastest
third-party implementation available. Two scripts drive everything. One
regenerates \papertable{} (time and peak heap of each form, and the fastest
baseline). The other regenerates \paperablation{}, which adds the scheduling
optimizations one at a time. The apps check every form against a reference
before timing it.

\subsection{Artifact check-list (meta-information)}

{\small
\begin{itemize}
  \item {\bf Algorithm:} inductive Funcs in Halide: a recurrence defined over
    its own earlier values, lowered with fusion, sliding window and storage
    folding.
  \item {\bf Program:} 9 benchmark apps: random number generation, a biquad
    filter cascade, pairwise sequence alignment, Viterbi decoding, an ODE
    integrator (RK4, Allen--Cahn), prefix mean, Chebyshev iteration, and
    Mamba-2 and flash attention on the GPU.
  \item {\bf Compilation:} Halide built with LLVM~22 (installed by
    \texttt{uv}); C++17 with \texttt{-O3 -march=native}; clang for the biquad
    app; CUDA~13 for the GPU apps.
  \item {\bf Data set:} generated by each app (seeded); no downloads.
  \item {\bf Run-time environment:} Linux x86-64. jemalloc (\texttt{libjemalloc2})
    is required. Optional baselines: Boost.odeint, oneTBB, Intel IPP,
    parasail, Julia~1.12, scipy, torch, Triton, mamba\_ssm.
  \item {\bf Hardware:} an x86-64 CPU with AVX-512 and many cores (the paper
    used a 32-core workstation CPU). AVX2 works, at lower speed. An NVIDIA
    GPU with CUDA~13 for the three GPU rows.
  \item {\bf Execution:} two scripts; single-thread rows are pinned to one
    core, all-cores rows use one thread per physical core.
  \item {\bf Metrics:} best time over 30 trials (ms); peak Halide heap (MB)
    from the Halide profiler; resident memory growth for the ablation.
  \item {\bf Output:} Markdown tables and per-run logs under
    \texttt{inductive\_artifact/results/}.
  \item {\bf Experiments:} \papertable{} and \paperablation{}.
  \item {\bf How much disk space required (approximately)?:} 10~GB (LLVM,
    Halide build, Python baselines).
  \item {\bf How much time is needed to prepare workflow (approximately)?:}
    under an hour to build Halide; longer with the optional baselines.
  \item {\bf How much time is needed to complete experiments (approximately)?:}
    a few hours for both scripts; each app can be run alone.
  \item {\bf Publicly available?:} yes.
  \item {\bf Code licenses (if publicly available)?:} MIT (Halide); the
    vendored baselines keep their own licenses.
  \item {\bf Archived (provide DOI)?:} \artifactdoi.
\end{itemize}
}

\subsection{Description}

\subsubsection{How delivered}
The branch \texttt{inductive\_artifact} of \artifactrepo, archived at
\artifactdoi. The benchmarks are in \texttt{apps/}; the artifact's scripts are
in \texttt{inductive\_artifact/}.

\subsubsection{Hardware dependencies}
The CPU experiments need an x86-64 CPU. The paper's machine is a 32-core AVX-512
workstation CPU (four 8-core chiplets, 128~MB L3, one NUMA node). The
all-cores rows are memory-bound, so their ratios depend on the machine's memory
bandwidth per core. The GPU rows need an NVIDIA GPU with tensor cores and
CUDA~13; the paper used a consumer GPU.

\subsubsection{Software dependencies}
Linux, a C++17 compiler, GNU make, CMake, \texttt{uv}, and jemalloc. The driver
preloads jemalloc into every process, configured to keep freed memory, so that no
timed run pays page faults for scratch that an earlier run freed. Optional
baselines are built when found: Boost.odeint (\texttt{BOOST\_DIR}, or
\texttt{libboost-dev}), oneTBB (\texttt{TBB\_DIR}, or \texttt{libtbb-dev}), Intel
IPP (\texttt{IPP\_DIR}), parasail (\texttt{PARASAIL\_DIR}), Julia
(\texttt{HB\_JULIA}, or \texttt{julia} on the \texttt{PATH}), and the Python
packages in \texttt{inductive\_artifact/requirements.txt}. A row whose
baseline is missing reports the fastest baseline that is present.

\subsubsection{Data sets}
None; every app generates its input from a fixed seed.

\subsection{Installation}

{\footnotesize
\begin{verbatim}
git clone -b inductive_artifact \
  https://github.com/wraith1995/Halide.git
cd Halide
sudo apt install libjemalloc2 libboost-dev \
  libtbb-dev
inductive_artifact/setup.sh
# or, with the Python and GPU baselines:
WITH_PYTHON_BASELINES=1 \
  inductive_artifact/setup.sh
\end{verbatim}
}

\texttt{setup.sh} installs LLVM into \texttt{.venv} with \texttt{uv} and builds
the Halide distribution in \texttt{distrib/}, which every app links against.

\subsection{Experiment workflow}

{\footnotesize
\begin{verbatim}
inductive_artifact/run_benchmarks.sh  # table
inductive_artifact/run_ablation.sh    # ablation
\end{verbatim}
}

Both scripts take \texttt{--only app,app} to run a subset (app names:
\texttt{cpu\_rng}, \texttt{cpu\_biquads}, \texttt{cpu\_alignment},
\texttt{viterbi}, \texttt{ode}, \texttt{prefixsum}, \texttt{chebyshev},
\texttt{cuda\_mamba2}, \texttt{flash\_attention}). Rows are cached in
\texttt{apps/inductive\_benchmarks\_logs/}, so a partial run updates the
table instead of replacing it. The GPU ablation levels are printed by
the GPU apps themselves:
{\footnotesize
\begin{verbatim}
cd apps/cuda_mamba2
make SCAN=inductive WMMA=true \
  ABLATE_FUSION=1 test test_bwd
cd ../cuda_attention && make test
\end{verbatim}
}
Their \texttt{ablation fusion}, \texttt{ablation slide} and
\texttt{ablation slide+unroll} rows give the GPU ablation entries.

\subsection{Evaluation and expected results}

The scripts write \path{benchmarks.md}, \path{ablation_time.txt},
\path{ablation_memory.txt} and \path{ablation_vs_thirdparty.txt} in
\path{inductive_artifact/results/}. The paper machine's outputs are in
\path{apps/inductive_benchmarks_results.md} and
\path{apps/inductive_ablation_results/} for comparison. The claims to check:

\begin{enumerate}
  \item \textbf{Correctness.} Every form passes its app's check against the
    reference before it is timed; a failed check stops the row.
  \item \textbf{Memory.} The inductive form's peak heap is zero, or only the
    part that cannot fold (alignment's traceback plane). The RDom form's peak
    heap is the whole trajectory, e.g.\ 2147~MB for biquads and 1074~MB for
    the all-cores ODE row. These numbers are properties of the program and
    should match exactly on any machine.
  \item \textbf{Time, all cores.} The inductive form beats the RDom form where
    many recurrences share the machine's memory bandwidth. On the paper
    machine the margin is 2--19$\times$ (prefix mean 2.0$\times$, Viterbi
    4.4$\times$, rng 5.9$\times$, biquads 7.4$\times$, alignment
    16.5$\times$, ODE 19$\times$). The inductive form is at most 7\% slower
    than the fastest third-party baseline (biquads), and faster than it for
    alignment and ODE.
  \item \textbf{GPU.} The inductive form is 4--15\% faster than its RDom
    form (Mamba-2 forward and backward, flash attention). It is within 6\% of
    the vendor kernels for the Mamba-2 forward and flash attention; Triton's
    Mamba-2 backward is 1.4$\times$ faster.
  \item \textbf{Controls.} Chebyshev and the single-thread Viterbi rows take
    about the same time in both forms (within 20\%): their trajectories fit in
    cache, or each step is a long dependency chain compared with the bytes it
    writes.
  \item \textbf{Ablation.} Fusion alone is far slower than every other level,
    since each output recomputes the recurrence from its start (62$\times$ to
    6038$\times$ the paper form on the CPU, measured at a small size).
    Adding the sliding window removes the recomputation and lands near the
    RDom form (1.3--14$\times$ the paper form), except for ODE, where the
    unfolded window holds three times the RDom form's storage (39$\times$).
    Storage folding then gives the paper form, the fastest level in every
    experiment.
\end{enumerate}

Absolute times depend on the machine. The all-cores ratios grow with the
number of cores sharing each memory controller, so a machine with fewer cores
or more bandwidth per core shows smaller all-cores ratios; their order across
apps, and the memory columns, should not change.

\subsection{Experiment customization}

Each app's Makefile takes its problem sizes on the command line (e.g.\
\texttt{make -C apps/cpu\_biquads SECTIONS=8 CHANNELS=1024 SAMPLES=262144
PAR=true test}); the README in each app lists them. The driver's environment
variables set the core count (\texttt{HB\_NCORES}), the core that
single-thread rows are pinned to (\texttt{HB\_PIN\_CORE}), the Python used for
the baselines (\texttt{HB\_PYTHON}), and the jemalloc library
(\texttt{HB\_JEMALLOC}).

\subsection{Notes}

The ODE app schedules its output's zero fill across the batch
(\texttt{E.reorder(n, b).vectorize(n, 16).parallel(b)}). Without that line
the fill runs serially in every form and adds about 1.3~ms to each run.
\texttt{apps/inductive\_ablation\_results/README.md} describes this fix and
gives the paper-commit source with it applied
(\texttt{ode\_paper\_close/}).

\subsection{Methodology}

Submission, reviewing and badging methodology:
\begin{itemize}
  \item \url{https://www.acm.org/publications/policies/artifact-review-and-badging-current}
  \item \url{https://cTuning.org/ae}
\end{itemize}
